\section{One loss and one threshold}\label{sec:epsilon}

Fix the rational vector in \eqref{eq:old:8_1} and the complete analytic loss
\(\beta_0=g_0/16\) from \eqref{eq:old:8_4}. In every determinant application, the positive exponent gap is fixed
first. The degree and the matrix dimensions are chosen after it. This
gives a finite family of degrees and dimensions. A moving box may select
one member of it. The coefficients, the polynomial and the resultant may
vary with the box and with \(Y\). Only their degrees and height exponents
are uniform. Nonvanishing comes from the source determinant, from
coprimality, or from the endpoint alternative in the relevant row. The
height bound plays no part in it.

Fix a selected member. The cofactor solution and the explicit
integerizer give a polynomial height \(Y^{C_i}\). Here \(C_i\) depends
only on the fixed degree, the dimension and the allocated loss. Hence a
nonzero resultant is an integer of size at most \(Y^{C_i'}\)
for a fixed exponent \(C_i'\). The integer itself may change as the
source labels move.

We allocate the loss as follows. Suppose one complete row proof contains
\(N_i\) divisor, Heath--Brown, square-fibre, or lattice estimates. Let
the \(j\)-th estimate apply to a parameter bounded by \(Y^{h_{ij}}\).
Here \(N_i\) is the fixed finite number of estimate stages in the complete
row proof. A stage repeated over moving boxes keeps its one uniform
exponent choice. It does not incur a new loss for each box.
Choose \(h_{ij}\) at least as large as every polynomial-height exponent in
that invocation. This covers the auxiliary form, its integerization, any
nonzero resultant and the size exponent of the counted parameter.
For a product of height factors, add their exponents when choosing
\(h_{ij}\). Then it bounds the total contribution of the product.
The degrees and height exponents are uniform over the complete moving
source family in that row.
Choose the internal exponent \(\rho_{ij}\) so that
\[
0<\rho_{ij}\le
\frac{\beta_0}{8N_i(1+h_{ij})}.
\AutoTag{eq:old:9_1}
\]
Then
\[
\sum_{j=1}^{N_i}h_{ij}\rho_{ij}\le\frac{\beta_0}{8}.
\]
Allocate the other fixed portions of \(\beta_0\) to the finitely many
dyadic sums, logarithms and divisor multiplicities in the same complete
row proof. This is one allocation for the whole row. After a row
estimate is stated, no factor or inverse count is charged again.

The thirteen first-match estimates are stated in
Propositions~\ref{prop:outerrows} and~\ref{prop:q2}, and in the row
Propositions~\ref{prop:row-3u}, \ref{prop:row-3t},
\ref{prop:row-4e}, \ref{prop:row-4h}, \ref{prop:row-4p}, and
\ref{prop:row-5w}. They use the finite cutoff data summarized below. Each cutoff is fixed
by the corresponding proof and depends only on \(\beta_0\) and the
rational vector. The table records conditions and the estimates that
supply them; it does not give numerical values for all row thresholds
or analytic implied constants. Equation \eqref{eq:old:9_5} takes the
maximum of the thresholds used in those proofs, including all displayed
row conditions.

\begin{table}[!ht]
\centering
\caption{Threshold conditions for the thirteen row estimates.}
\label{tab:threshold-conditions}
\begin{tabular}{@{}>{\raggedright\arraybackslash}p{0.14\textwidth}>{\raggedright\arraybackslash}p{0.79\textwidth}@{}}
Row & Threshold conditions included in $Y_i(\beta_0)$\\ \hline
$O_R$ & $Y\ge64$, $H=Y/2\ge2$, $V=Y^{5/3-\epsilon_0}<H^2$; fixed-height partner and low-radical estimates\\
$O_M$ & Ordered residual-factor divisor bound and the fixed-height quadratic-fibre estimate\\
$O_{\operatorname{rad}M}$ & One fixed positive Rankin parameter and the same quadratic-fibre estimate\\
$Q0$ & $Y\ge256$ and $3Y^{q_{\max}-1}<1$\\
$Q1$ & Fixed-height bound for all $(P,L)$ labels\\
$Q2u$ & $1/3-5\epsilon_0>0$; unit-root count, one $u$-sum, and the dyadic sector count\\
$Q2t$ & $1/3-6\epsilon_0-l_0>0$; corresponding $t$-chart count and one $t$-sum\\
$3u$ & Positive determinant gap, finite factor family, all-aspect lattice count, axes, and inverse bounds\\
$3t$ & Determinant cutoff, fixed grid condition $N\ge40$, and
later lattice, factor, census, and inverse cutoffs\\
$4E$ & Four supported-factor sums, their fixed-height bounds, and the single shift sum\\
$4H$ & Reducible, content, and primitive cubic branches, including the primitive-form estimate\\
$4P$ & Graph, high-degree, cover, and axis bounds, with simultaneous shifts counted once\\
$5W$ & Lemma~\ref{lem:ratio-cell-determinant}, finite chart comparisons,
original-grid census, factor branches, endpoint bounds, assigned-corner
cases, and the geometric envelope; in particular $6E_Y<\xi/2$ and
$576Y^{q_{\max}+2\lambda-\tau_0}<1$ in \eqref{eq:old:7_15}\\
\end{tabular}
\end{table}

\begin{samepage}
In the \(5W\) row, the geometric cutoff is
\[
T_{\rm geom}=\frac{252}{g_0}+1.
\AutoTag{eq:old:9_3}
\]
\end{samepage}%
For \(\log Y>T_{\rm geom}\), \(\log2<7/10\) gives
\[
\frac{90\log2}{\log Y}
<\frac{63}{252/g_0}=\frac{g_0}{4}.
\AutoTag{eq:old:9_4}
\]
The high-degree correction in \eqref{eq:old:7_25} is separate from the geometric
correction. Since \(a_{\rm grid}>1/2\), for \(Y\ge2\),
\[
\lambda+\frac{c_N}{2}
\le \frac{Y^{-1}+\tfrac12Y^{-1/2}}{\log Y}
<\frac1{\log Y}.
\]
Also, \(c_N=\log_Y(1+Y^{-a_{\rm grid}})
\le Y^{-a_{\rm grid}}/\log Y\), so the interval correction is smaller
than \(1/\log Y\).
Consequently
\[
\frac{E_q+c_N}{2}+12d_2
\le\frac{\tfrac12\log18+12\log2+1}{\log Y}
<\frac{90\log2}{\log Y}<\frac{g_0}{4}.
\]
Here \(\tfrac12\log18+12\log2=9.763\ldots\). Thus each of the geometric
correction \(6\rho_Y\), the interval correction \(c_N\), and the high-degree
correction \((E_q+c_N)/2+12d_2\) is at most \(g_0/4\). The final row estimate
is charged their maximum, not their sum.
We retain this conservative charge in the reserve \eqref{eq:old:8_4};
it is unnecessary for the exponent in Proposition~\ref{prop:row-5w},
where the three finite corrections are absorbed into the implied constant.
The determinant, census and geometric-envelope thresholds in that row
are the conditions used in Section~\ref{sec:long}. The
contribution \(90\) is the source-mass \(5\) plus the one geometric
collar mass \(1\), each times the coefficient \(15\) in the
row margin.

Let \(Y_i(\beta_0)\) be the maximum of the finitely many thresholds
used in the proof of row \(i\), after applying \eqref{eq:old:9_1}.
Define
\[
T_0=\max\{3,T_{\rm geom},\log Y_1(\beta_0),\ldots,
                         \log Y_{13}(\beta_0)\},\qquad
Y_{\rm final}=\lceil e^{T_0}\rceil+1.
\AutoTag{eq:old:9_5}
\]
Since \(T_0\ge3\), \(Y_{\rm final}\ge\lceil e^3\rceil+1=22>4\).
Thus the elementary conditions \(Y\ge4\) of the 5W accounting
hold at the common cutoff.
The maximum runs over the thirteen fixed row estimates and their fixed
finite degree families. It does not run over moving source labels. It
includes the separate post-determinant thresholds of the \(3t\) row.

\begin{proposition}[Common-loss assembly]\label{prop:commonloss}
For \(Y\ge Y_{\rm final}\), each row \(i\ne5W\) is
\[
\#\mathrm{row}_i\ll_{\beta_0}
Y^{2/3-r_i+\beta_0},
\qquad r_i\ge S_0,
\]
and the last row satisfies
\[
\#5W\ll_{\beta_0}
Y^{2/3-S_0+\beta_0+g_0/4}.
\]
Consequently \(\#\Omega_Y\ll_\varepsilon
Y^{109/165+\varepsilon}\).
\end{proposition}

\begin{proof}
The first twelve row estimates are complete counts. They include the
moving labels, shifts, coefficient factorizations and inverse
multiplicities. By Section~\ref{sec:ledger}, their minimum saving is
\(S_0\). By \eqref{eq:old:7_26}, the three limiting values
\(L_{5W},s_{\rm int},s_{\rm HD}\) in the definition of \(\sigma_{5W}\)
are each at least \(S_0\). Hence the corrected minimum satisfies
\[
 \sigma_{5W}\ge S_0-\max\!\left(6\rho_Y,\ c_N,\
       \frac{E_q+c_N}{2}+12d_2\right).
\]
The three corresponding source-box contributions are added in
\eqref{eq:old:7_25}, so this minimum controls their sum exponent.
The corrections have separate sources. First, \(6\rho_Y\) is the finite
source-to-lower-corner allowance in \eqref{eq:old:7_22}--\eqref{eq:old:7_23c}.
Next, \(c_N\) is the ceiling correction from \(N=\lceil Y^{a_{\rm grid}}\rceil\)
in \eqref{eq:old:7_6}. Last, \((E_q+c_N)/2+12d_2\) bounds the finite
square-root term \(\sqrt{4MN}\). This uses the source-product and
grid-ceiling bounds just before \eqref{eq:old:7_25}. The margins add no
losses beyond these three terms.
By \eqref{eq:old:9_4} and the correction bounds above, each of the three
corrections is at most \(g_0/4\) for \(Y\ge Y_{\rm final}\). Hence
\(\sigma_{5W}\ge S_0-g_0/4\). Apply \eqref{eq:old:7_25} with
\(\eta=\beta_0\), using the remaining \(\beta_0/2\) for the box sum.
Thus the total count is
\[
\ll_{\beta_0}
Y^{2/3-S_0+\beta_0+g_0/4}
=Y^{109/165-11g_0/16}\ll Y^{109/165}.
\]
There are thirteen rows. Their finite sum changes only the implied
constant.
\end{proof}
